\documentclass[12pt]{article}
\usepackage{amsmath,amssymb}
\begin{document}
Scratchpad for the rectangle 2-dilation problem.

Interpretation: rectangles are oriented Euclidean boxes; a map
\(f:(R,\partial R)\to(S,\partial S)\) of degree 1 means the induced map on
\(H_4(\cdot,\partial\cdot;\mathbb Z)\) sends \([R]\) to \([S]\).
The constant called \(k\) in the problem is a small universal constant;
in the solution it is denoted \(\kappa\) to avoid conflict with
\(k\)-dilation.

Main input used in answer.tex: the \(n=4,k=2\) two-parameter
Guth tightening/cup-square estimate.  A degree-one two-parameter sweepout by relative 2-cycles satisfies
\[
 \frac{M_2}{S_3S_4}+
 \frac{M_3}{S_1S_3S_4}+
 \frac{M_4}{S_1S_2^{1/2}S_3^{3/2}S_4}\ge c_0,
\]
where \(M_2\) bounds 2-cycle masses, \(M_3\) bounds the 3-chain swept out by moving across the first parameter direction, and \(M_4\) is the total 4-chain mass.  This is the specialized quantitative
sweepout estimate from Guth's thesis / area-contracting maps between
rectangles, proved by the Federer--Fleming/Guth tightening procedure.

Application: slice \(R\) by the first two coordinates.  Vertices are
\(x_3x_4\)-slices, first-direction edges are \(x_1x_3x_4\)-slabs, and
squares are four-dimensional blocks.  Since \(\Dil_2(f)\le 1\), singular
values show also \(\Dil_3(f),\Dil_4(f)\le1\).  Hence
\(M_2\le R_3R_4\), \(M_3\le R_1R_3R_4\), \(M_4\le V(R)\).
Substitution gives
\[
 R_3R_4/(S_3S_4)+R_1R_3R_4/(S_1S_3S_4)
 +V(R)/(S_1S_2^{1/2}S_3^{3/2}S_4)\ge c_0.
\]
Choosing \(\kappa=\min(1/10,c_0/4)\) proves the dichotomy.
\end{document}
